\section{考察}
\subsection{3つの構成の使い分け}
COCOの80クラスの範囲では，深さ0の精度はYOLO11xが提案手法を大きく上回り，その差は検出器ではなくカテゴリの分類器に起因する．
ハイブリッド手法は，深さ0をYOLO11xに任せることでこの差を解消し，粒度を変更できたインスタンスを提案手法の1.57〜1.79倍に増やした．
したがって，深さ0のカテゴリが閉域の検出器の語彙に含まれる場合はハイブリッド手法が適する．
一方，提案手法はDETRのボックスに対して任意の語彙から深さ0を選べるため，深さ0の語彙自体をCOCOから変更したい場合に用いることができる．
ただし，DETRもCOCOで学習されているため，提案手法でもCOCO外の物体の検出は11.0\%にとどまり，語彙を自由にできることがCOCO外の物体への対応に直結しなかった．

\subsection{COCO外の物体への対応}
COCO外の物体の粒度の変更は，検出，深さ0の命名，下位語の選択の3段階の積で決まる．
検出の段階では，COCOで学習した検出器の閾値を下げても再現率は2割未満で飽和し，見落とした物体は後段で回復できない．
物体らしさを学習したOWLv2の候補をYOLO11xのボックスと重ならない領域に限って加えることで，COCOの物体の精度を保ったまま検出を6.9倍にでき，検出器が粒度の変更の上限を決めることが確認された．
一方，命名の段階では，検出したCOCO外の物体のうち深さ0を正しく命名できたのは8.1\%であった．
LVISのカテゴリ名を語彙とした場合でも26.7\%であるため，命名の誤りは語彙の不足だけでなく，SigLIP-ViT-B/16による小さな領域の分類の難しさにも起因する．
WordNetの同じ深さの語彙には，物体の部分（leg，screen）や多義の語（coating，slit）が含まれ，加えた候補の約8割はGTと対応しない領域であったため，これらの領域に不適切な語が与えられた．
したがって，物体の名前としての用いられやすさを考慮した語彙の選択と，確信度による領域の棄却が必要である．
\add{4.6節の結果から，物体らしさのスコアによる棄却は加えた候補の適合率を高め，粗い語から順に選ぶ命名は，VLMの能力が高い場合に，最初から細かい語を1つ選ぶ命名を大きく上回った．一方，キャプションに出現する語への語彙の限定は，最初から1語を選ぶ命名では効果が無く，語彙の大きさよりも，選ぶ順序と停止の仕組みが重要であった．ただし，改善後もCOCO外の物体の命名は4.1\%にとどまり，検出した物体の多くは正しく命名できていない．}

\subsection{今後の課題}
第一に，深さ0の語彙を，画像データセットにおける出現頻度や語の具体性に基づいて選び直すことが挙げられる．
第二に，\add{4.6節で評価した棄却は，加える候補の数の予算に対してCOCO外の物体の再現率を最大化する閾値を選んだものであり，LVISの注釈が網羅的でないため加えた候補の適合率を正確に評価できていない．網羅的な注釈を持つデータセットでの評価と，}加えた候補のうち物体でない領域や物体の部分を，シグモイドの確信度やOWLv2の物体らしさのスコアにより棄却する規則の設計が挙げられる．
第三に，シグモイドは誤った枝の検出に，ソフトマックスは兄弟の語の間の比較に適するため，両者を組み合わせた停止規則が挙げられる．
第四に，より規模の大きいVLMや，LVIS以外の細かい粒度のアノテーションを持つデータセットでの評価が挙げられる．
